% Source PDF: source/普通高考/2013/2013新课标2文(贵州,甘肃,青海,西藏,黑龙江,吉林,海南,宁夏,内蒙古,新疆,云南).pdf, page 1, question 7.
% PDF embedded bitmap attempt: pdfimages -list returned no image objects; the flowchart is vector page content.
% Grid evidence: tmp/review_2013_ncp2_liberal/grid/page1_grid100.jpg and q7_grid2.jpg.
% Original-side comparison crop: tmp/review_2013_ncp2_liberal/crops/q7_source_tight.png from a no-grid page render.
% Elements: rounded start/end terminals, input/output parallelograms, four process
% rectangles, one decision diamond, directed connectors, loop-back branch, 是/否 labels.
% Node-edge list: start->input N->(k=1,S=0,T=1)->(T=T/k)->(S=S+T)->(k=k+1)
% ->(k>N); (k>N)-是->输出S->结束; (k>N)-否->loop-back to the stem before T=T/k.
% No loop edge enters the output branch; all connectors are solid and directed.
% solid segments: every node border and connector/loop segment; dashed segments: none.
% Labels: 是 is beside the downward output branch, 否 is above the right loop branch;
% all node text is centered with local zero indentation and no manual horizontal offset.
\begin{tikzpicture}[
  >=Stealth,
  x=1cm,y=1cm,
  line width=0.45pt,
  line cap=round,line join=round,
  font=\normalsize,
  flowtext/.style={anchor=center,align=center,inner sep=0pt,
    execute at begin node={\setlength{\parindent}{0pt}}},
  flowbox/.style={flowtext,draw,minimum height=.68cm,inner xsep=4pt,inner ysep=2pt},
  terminal/.style={flowbox,rounded corners=.16cm,minimum width=1.25cm},
  process/.style={flowbox},
  io/.style={flowbox,shape=trapezium,trapezium left angle=72,trapezium right angle=108,
    minimum height=.75cm},
  decision/.style={flowbox,diamond,aspect=2.8,minimum width=3.45cm,minimum height=.90cm,
    inner sep=0pt}
]
  % Keep a small safety margin below the rounded end terminal; do not clip its border.
  \path[use as bounding box] (-1.85,-10.05) rectangle (2.30,0.45);

  \node[terminal] (start) at (0,0) {开始};
  \node[io,minimum width=1.90cm] (input) at (0,-1.10) {输入 $N$};
  \node[process,minimum width=3.80cm] (init) at (0,-2.28) {$k=1,S=0,T=1$};
  \node[process,minimum width=1.80cm,minimum height=.90cm] (divide) at (0,-3.75)
    {$T=\dfrac{T}{k}$};
  \node[process,minimum width=2.30cm] (accumulate) at (0,-4.92) {$S=S+T$};
  \node[process,minimum width=1.90cm] (increment) at (0,-6.12) {$k=k+1$};
  \node[decision] (test) at (0,-7.35) {$k>N$};
  \node[io,minimum width=1.80cm] (output) at (0,-8.45) {输出 $S$};
  \node[terminal] (finish) at (0,-9.58) {结束};

  % Main chain and the true/output branch.
  \draw[->] (start.south) -- (input.north);
  \draw[->] (input.south) -- (init.north);
  \draw[->] (divide.south) -- (accumulate.north);
  \draw[->] (accumulate.south) -- (increment.north);
  \draw[->] (increment.south) -- (test.north);
  \draw[->] (test.south) -- node[right=2pt] {是} (output.north);
  \draw[->] (output.south) -- (finish.north);

  % False branch: right, up, then left-arrow merge into the main stem before T=T/k.
  \coordinate (merge-level) at (0,-2.95);
  \coordinate (loop-right) at (2.20,-7.35);
  \draw (test.east) -- (loop-right) -- (2.20,-2.95);
  \draw[->] (2.20,-2.95) -- (0.40,-2.95) -- (merge-level);
  % The central stem carries the sole downward arrow into the divide node.
  \draw (init.south) -- (merge-level);
  \draw[->] (merge-level) -- (divide.north);
  \node[flowtext,anchor=west,inner sep=1pt] at (1.48,-7.05) {否};
\end{tikzpicture}
